\documentclass{article}
\usepackage{graphicx}
\usepackage{booktabs}
\usepackage{multirow}
\usepackage{array}
\usepackage{amsmath}
\usepackage{amssymb}
\usepackage{wrapfig}
\usepackage{adjustbox}

\usepackage{corl_2026} % Use this for the initial submission.
% \usepackage[final]{corl_2026} % Uncomment for the camera-ready ``final'' version.
% \usepackage[preprint]{corl_2026} % Uncomment for pre-prints (e.g., arxiv); This is like ``final'', but will remove the CORL footnote.

\title{Code-based Real-world Automated Data Acquisition via Information-guided Trajectory Selection}
% \title{Code-based Automated Data Acquisition for Real-world Robot Learning}

% The \author macro works with any number of authors. There are two
% commands used to separate the names and addresses of multiple
% authors: \And and \AND.
%
% Using \And between authors leaves it to LaTeX to determine where to
% break the lines. Using \AND forces a line break at that point. So,
% if LaTeX puts 3 of 4 authors names on the first line, and the last
% on the second line, try using \AND instead of \And before the third
% author name.

% NOTE: authors will be visible only in the camera-ready and preprint versions (i.e., when using the option 'final' or 'preprint'). 
% 	For the initial submission the authors will be anonymized.

\author{
  Jane E.~Doe\\
  Department of Electrical Engineering and Computer Sciences\\
  University of California Berkeley 
  United States\\
  \texttt{janedoe@berkeley.edu} \\
  %% examples of more authors
  %% \And
  %% Coauthor \\
  %% Affiliation \\
  %% Address \\
  %% \texttt{email} \\
  %% \AND
  %% Coauthor \\
  %% Affiliation \\
  %% Address \\
  %% \texttt{email} \\
  %% \And
  %% Coauthor \\
  %% Affiliation \\
  %% Address \\
  %% \texttt{email} \\
  %% \And
  %% Coauthor \\
  %% Affiliation \\
  %% Address \\
  %% \texttt{email} \\
}


\begin{document}
\maketitle

%===============================================================================

% Abstract <초안 1>
    % The advent of Vision-Language-Action (VLA) models and the paradigm shift toward data-driven robot learning have intensified the demand for high-quality, large-scale robotic datasets. However, current data acquisition methodologies remain predominantly reliant on human teleoperation, which presents two critical limitations. First, this manual dependency creates a severe scalability bottleneck; adapting to novel domains, environments, or embodiments necessitates extensive data recollection, imposing prohibitive human effort. Second, teleoperated data suffers from sparse semantic labeling, where high-level instructions are typically assigned only at the episode level. This lack of granularity constrains policies to act as mere trajectory followers without any semantic understanding, rendering the learning process increasingly difficult as task horizons expand. To address these challenges, we propose the Robo-Scrape framework, an automated data acquisition pipeline that leverages CodeLLMs to generate skill-level dense semantic labels and task trajectories. By automating the collection process and providing fine-grained semantic annotations, our framework improves both data scalability and policy learnability. To validate our approach with real-world data, we curated a benchmark covering eight common robotic manipulation tasks. Experimental results demonstrate that policies trained on data generated by our framework achieve performance comparable to those trained on conventional human-teleoperated data. All code, datasets, and models are publicly available at [Project URL].

% Abstract <초안2>
 % 초안2
    % 1. 로봇 연구의 패러다임이 바뀌면서 높은 퀄리티와 대량의 데이터 취득이 중요해지고 있어. 따라서 확장성이 부족한 사람의 텔레옵을 대체하기 위한 다양한 연구들이 진행되고 있어. 
    
    % 2. 그 중 하나는 Code based policy를 활용한 데이터 자동 취득 연구야. 해당 연구들은 시뮬레이션 환경에서 확장가능한 데이터취득 자동화를 구축지만, 실세계에 곧바로 적용하기엔 아래 3가지 문제가 있어. (1) absence of simulation privileges (2) real-world transition variability (3) embodiment-bound acquisition cost
    
    % 3. 위 문제들을 해결하기 위해서 CRAPE를 제안함. CRAPE는 실세계 데이터취득을 고려한 자동화 기능을 제공함과 pre-selective 데이터취득으로 높은 데이터 퀄리티를 보장하는 프레임워크야. 
    
    % 4. By considering the 완전한 무인화, 테스크 수행 및 초기화, scene layout randomization을 스케쥴링하여 자동화하였고. 기존 code based policy의 문제인 유사한 데이터 취득을 해결하기 위한 다양화 방법, 효율적 데이터취득 및 학습을 위한 pre-selective acquisition 방법을 제공함...
% <초안2> 반영 - 26.05.03
\begin{abstract}
   Vision-Language-Action (VLA) models provide a unified framework for general purpose robotic manipulation, but their performance is still limited by costly human teleoperation. Executable policy-based data generation, such as Code as Policies (CaP), offers a scalable alternative, yet its deterministic primitives often produce redundant trajectories and action ambiguity that destabilize flow matching or diffusion based VLAs. While Mutual Information (MI) is useful for offline curation, it is unreliable for online cold start acquisition, where local action statistics are sparse and selection remains trapped in narrow warm start coverage. This motivates separating coverage expansion from action distribution refinement. We propose CRAPE, a Code-based Real-world autonomous data Acquisition PipEline that decouples state coverage formation from action distribution refinement. CRAPE first seeds state coverage via canonical subgoal exploration, expanding the workspace support while installing clean local action references. Within that support, it then scores each automatically generated candidate \emph{before committing it to the VLA training set}, using two complementary axes: a buffer-side marginal MI-inspired coverage-consistency score gain, and a VLA-side informativeness term computed by a denoising-loss evaluation of VLA. On a 10-task SO-101~\citep{cadene2024lerobot} real-world benchmark, CRAPE outperforms unfiltered code-based collection and reaches teleoperation-level macro-average performance across VLA backbones. Ablations show that state seeding provides the largest gain, while pre-selective admission converts raw trajectory diversity into learnable supervision.
\end{abstract}
% \begin{abstract} 
%     The recent shift toward robotic foundation models has made large-scale, high-quality robot data a central requirement for robot learning. However, conventional human teleoperation is inherently difficult to scale. As an alternative, code-based policy approaches have been proposed to automate data acquisition in simulation. These methods demonstrate the potential of scalable data generation, but their assumptions are largely tied to simulation. Directly applying these works to real-world data acquisition poses three key challenges: absence of simulation privileges, real-world transition variability, and embodiment-bound acquisition cost. First, real-world environments do not provide task-relevant information, workspace reset, and randomization, all of which are readily available in simulation effortlessly. Second, perception noise and hardware jitter induce variable transitions that frequently push the robot into OOD states, whereas conventional code-based policies provide limited state coverage. Third, since each rollout occupies the physical embodiment, redundant and useless data collection without pre-selective criteria leads to acquisition inefficiency. To address these limitations, we propose CRAPE, a real-world code-based robot data acquisition framework that enables automated, diverse, and pre-selective data collection.
%     We evaluate CRAPE on real-world robotic manipulation tasks ... Experimental results show that ...
% \end{abstract}

%===============================================================================
% Two or three meaningful keywords should be added here
% Real-World Data Acquisition, Vision-Language-Action Models, Code-as-Policies, Robot Manipulation
\keywords{Robotic data acquisition, Robot foundation models, Vision-Language-Action Models} 

%===============================================================================
\section{Introduction}
    % ---------------- P1: Motivation — VLA + data dependency + automation frontier ----------------
    Vision-Language-Action (VLA) models~\citep{kim2024openvla, black2024pi_0, intelligence2025pi_, zitkovich2023rt, bjorck2025gr00t, team2024octo} have emerged as a unified paradigm for general-purpose robotic manipulation, directly mapping language instructions and visual observations to actions and exhibiting substantial generalization across objects, scenes, and embodiments. However, their capability is fundamentally bounded by the scale and diversity of training demonstrations~\citep{o2024open, khazatsky2024droid, walke2023bridgedata, wu2024robomind}, which are predominantly collected through human teleoperation, an approach that incurs prohibitive costs in labor, hardware wear, and embodiment coverage. Automated data collection has therefore become an important research direction, and Code-as-Policies (CaP)~\citep{liang2023code}, which composes skills into executable task-specific code, has emerged as a leading framework for both real-world manipulation~\citep{singh2023progprompt, huang2022inner, duan2024manipulate} and large-scale demonstration synthesis in simulation~\citep{wang2023robogen, mu2024robotwin}.

    % ---------------- P2: low diversity + incompatible-mode coexistence → why pre-selective; MI cold-start ----------------
    However, generating candidates at scale does not by itself produce a useful VLA training set. 
    Because CaP composes deterministic primitives along fixed waypoints, repeated rollouts are often redundant, while the variation that does appear can introduce action ambiguity: incompatible action modes under similar states, which destabilizes flow-matching and diffusion-based VLAs during training and closed-loop execution~\citep{black2024pi_0, bjorck2025gr00t}. 
    Thus, the key acquisition problem is not how many candidates are generated, but which ones provide useful supervision. 
    % 0529 박은영: 인용 논문 수정 cite{belkhale2023data} -> \citep{hejna2025robot} (mutual information을 IL 데이터 품질 기준으로 삼은 건 hejna2025robot 임)
    Mutual information between states and actions offers a natural criterion for this decision in imitation learning~\citep{hejna2025robot}, favoring action-coverage expansion without increasing conditional ambiguity. 
    Yet in online acquisition, MI is difficult to estimate in the cold-start regime, where the buffer contains too few local samples and selection can become trapped in the narrow coverage of the initial warm-start data.


\begin{figure*}[t]
  \centering
  \includegraphics[width=\textwidth]{figure/fig1.png}
    % 0526 kyung: left 랑 right만 있음 middle은 어디? 지금 left에 대한 설명은 telop 인데 설명에는 online MI-based라고 적혀있어서 불일치. 또한 middle의 경우도 어딜 봐서 이게 MI-based & cold start 라는게 드러나지 않음. 
    \caption{Overview of CRAPE. 
    % \textbf{(Left)} Naive online MI-based acquisition fails because the cold-start regime prevents reliable estimation of local action statistics, locking selection into the narrow coverage of warm-start trajectories. \textbf{(Right)} CRAPE decouples online acquisition into two phases. \emph{Phase~1} seeds state coverage via canonical subgoal exploration, expanding terminal-state regions while leaving each region with a clean canonical action seed. \emph{Phase~2} selects useful diversity within that support via marginal MI gain combined with VLA-side informativeness, prioritizing useful-OOD candidates over redundant-ID ones.
    }  % TODO: 최종 figure 확정 후 caption 다듬기
    \label{fig:introduction}
    \vspace{-0.1in}
\end{figure*}


    % ---------------- P3: CRAPE proposal — two-phase framework (no equations, split long sentences) ----------------
    We refer to this setting as a \emph{pre-selective acquisition problem}, in which candidates are selected online before VLA training rather than filtered after training on the full generated dataset.
    To address this problem, We propose CRAPE, a Code-based Real-world autonomous data Acquisition PipEline, as shown in Figure~\ref{fig:introduction}. 
    Phase~1 performs state-coverage seeding to diversify states and subgoals while keeping rollouts canonical, providing clean action seeds across the task support. 
    Phase~2 performs consistent diversity acquisition by combining a buffer-side marginal MI score, which expands action coverage without local ambiguity, with a VLA-side informativeness score from the Phase~1 model, which identifies candidates novel to the current policy. 
    CRAPE then selects the most VLA-informative candidate among the MI-useful ones, prioritizing useful OOD data while filtering out harmful OOD and redundant ID samples.
    % ---------------- P4: Evidence ----------------
    We evaluate CRAPE on the low-cost open-source SO-101 manipulator across real-world manipulation tasks and several VLA backbones. CRAPE-trained policies achieve higher performance comparable to human teleoperation, vanilla CaP and Manipulate-Anything~\citep{duan2024manipulate}.

    % ---------------- P5: Contributions ----------------
    Our contributions are threefold. 
    \textbf{(1)} We formulate code-based real-world data collection as online pre-selective acquisition and identify the cold-start limitation of directly applying MI-based curation during acquisition.
    \textbf{(2)} We propose CRAPE, a two-phase framework that first seeds broad state support through canonical subgoal exploration and then refines the action distribution through consistency-aware candidate admission.
    \textbf{(3)} We validate CRAPE on a 10-task SO-101 real-world benchmark with two VLA backbones, showing that selective code-generated data can outperform unfiltered code-based collection and reach teleoperation-level macro-average performance.
    % Our contributions are summarized as follows. \textbf{(1)} We formulate real-world automated data collection as an \emph{online pre-selective acquisition problem} and identify the cold-start limitation that prevents mutual-information-based curation from extending naturally to this setting. \textbf{(2)} We propose CRAPE, a two-phase acquisition framework that first seeds state coverage via canonical subgoal exploration and then refines the action distribution within that support via marginal MI gain. \textbf{(3)} We introduce an acquisition rule that distinguishes useful from harmful OOD by combining buffer-side MI usefulness with VLA-side informativeness, prioritizing candidates that are novel to the model yet consistent within the buffer. \textbf{(4)} We release an open-source SO-101 real-world benchmark, including skills, datasets, workspace specifications, and evaluation protocols, to facilitate reproducible VLA research.

%===============================================================================

\section{Related Work}
\textbf{Code-as-Policies and LLM-based robot data generation.}
LLMs and VLMs increasingly serve as programmatic interfaces for robot planning, manipulation, and demonstration synthesis. 
Code-as-Policies generates executable robot programs by composing perception APIs, geometric routines, and low-level primitives~\citep{liang2023code}, while ProgPrompt represents situated plans as structured programs over available objects and actions to improve executability~\citep{singh2023progprompt}. 
Beyond single-task execution, RoboGen and RoboTwin use code-based generation in simulation to synthesize large-scale robot demonstrations~\citep{wang2023robogen,chen2025robotwin}, and Manipulate-Anything extends VLM-driven code generation to real-world manipulation without demonstrations~\citep{duan2024manipulate}. 
These methods show that code can automate robot behavior and data generation, but they typically treat generated trajectories as demonstrations to be used directly. 
CRAPE instead treats them as candidates and asks which generated trajectories should be admitted into the VLA training set.

\textbf{Vision-Language-Action models.}
VLA models map visual observations and language instructions directly to robot actions, enabling language-conditioned policies to benefit from large-scale vision-language pretraining. 
RT-2 shows that such pretraining can transfer semantic knowledge to robotic control~\citep{zitkovich2023rt}, Octo trains an open-source generalist policy on cross-embodiment robot data for downstream adaptation~\citep{team2024octo}, and OpenVLA scales this direction with a 7B-parameter open-source VLA trained on real-world robot demonstrations~\citep{kim2024openvla}. 
As these backbones become more capable, real-world fine-tuning increasingly depends on the quality of the supervision provided to them. 
CRAPE is therefore complementary to VLA architecture design: it improves the data-acquisition process that determines what the backbone learns from.

\textbf{Data selection for imitation learning.}
Prior work shows that imitation learning depends not only on the amount of data but also on its quality. 
Mandlekar et al. analyze how demonstration quality affects downstream robot manipulation performance~\citep{mandlekar2021matters}, and Belkhale et al. study data quality through distribution shift, action divergence, and transition diversity~\citep{belkhale2023data}. 
More recently, Hejna et al. propose mutual-information-based estimators for robot data curation, selecting demonstrations that balance diversity and predictability in a fixed dataset~\citep{hejna2025robot}. 
These methods primarily address offline or post-hoc curation, where the dataset already exists. 
CRAPE studies the complementary online setting: candidates are scored during acquisition and committed only when they provide useful diversity without increasing local action ambiguity.
%===============================================================================
\section{Method}
\label{sec:method}
 
\subsection{Problem Formulation}
\label{sec:method:problem}

We consider an online data-acquisition pipeline for fine-tuning a VLA policy $\pi_\theta$ from a high-level natural-language task $I$. 
At acquisition step $t$, an automated code-based generator produces a candidate pool $\mathcal{C}_t = \{\xi^{(k)}\}_{k=1}^{K}$ over a skill sequence $z_{1:M}$, where each trajectory $\xi=(o_{1:T},u_{1:T})$ provides observation--action supervision for VLA training. 
Rather than adding all generated candidates to the training set, CRAPE scores the pool and commits only a selected subset $\mathcal{S}_t \subseteq \mathcal{C}_t$ to the VLA buffer,
\begin{equation} \small
\mathcal{D}_{t+1} = \mathcal{D}_t \cup \mathcal{S}_t .
\end{equation}
We call this setting \textbf{pre-selective acquisition problem}, since trajectories are selected during acquisition before final VLA training, rather than curated post hoc from a fixed dataset.

The goal is to retain \emph{useful diversity}: trajectories that expand action coverage without assigning incompatible actions to similar observations. 
Simply increasing the candidate pool is insufficient; redundant trajectories add little supervision, while locally inconsistent ones introduce \emph{action ambiguity} that can destabilize flow-matching and diffusion-based VLAs during training and closed-loop execution~\citep{black2024pi_0, bjorck2025gr00t}. 
% 0529 박은영: 인용 논문 수정 cite{belkhale2023data} -> \citep{hejna2025robot} (mutual information을 IL 데이터 품질 기준으로 삼은 건 hejna2025robot 임)
This coverage--ambiguity trade-off is captured by mutual information between states and actions in imitation-learning data~\citep{hejna2025robot}:
\begin{equation} \small
I_{\mathcal D}(S;A) \;=\; H_{\mathcal D}(A) \;-\; H_{\mathcal D}(A \mid S),
\label{eq:mi-classical}
\end{equation}
where $H_{\mathcal D}(A)$ rewards broad action coverage and $H_{\mathcal D}(A \mid S)$ penalizes conditional ambiguity. 
Since CRAPE builds the buffer online rather than scoring a fixed completed dataset, we evaluate each candidate $\xi$ by its marginal MI gain:
\begin{equation} \small
\mathrm{Gain}_{\mathrm{MI}}(\mathcal{D}_t, \xi)
\;=\;
I_{\mathcal{D}_t \cup \{\xi\}}(S;A)
-
I_{\mathcal{D}_t}(S;A)
\;=\;
\Delta H_A(\mathcal{D}_t, \xi)
-
\Delta H_{A \mid S}(\mathcal{D}_t, \xi).
\label{eq:marginal-mi-gain}
\end{equation}
Here, $\Delta H_A$ measures the action-coverage expansion introduced by $\xi$, whereas $\Delta H_{A \mid S}$ measures the increase in local action ambiguity around its observations. 
Thus, useful candidates should increase coverage while avoiding conflicting actions in already-covered state regions. 
However, this criterion requires sufficient state--action support to estimate local action statistics; in the cold-start regime, $\mathcal{D}_t$ is small and biased toward warm-start trajectories, making $\Delta H_{A \mid S}$ unreliable and motivating CRAPE's support-building phase before MI-based acquisition.

\subsection{Overall Framework}
\label{sec:method:framework}

To make the marginal MI criterion of Equation~\eqref{eq:marginal-mi-gain} usable in an online cold-start regime, we propose \textbf{CRAPE} (Figure~\ref{fig:introv}), a Code-based Real-world autonomous data Acquisition PipEline. The key design choice is to decouple state-support formation from action-distribution refinement, addressing the circular dependency that pure MI-based selection cannot escape on its own.

As illustrated in Figure~\ref{fig:introv}, \textbf{Phase 1} builds the clean support required for MI-based acquisition. 
Because reliable conditional action statistics cannot be estimated from a small buffer, CRAPE first constructs a seed buffer without using MI. 
It selects skill-wise subgoals that maximize terminal-state novelty and reaches each subgoal with a canonical interpolation planner, giving each visited region a single low-ambiguity action seed.
\textbf{Phase 2} refines this support through selective candidate admission.
Around each Phase-1 seed, CRAPE generates perturbed skill-level trajectories and scores them before adding them to the VLA training buffer.
The buffer-side score approximates marginal MI gain by rewarding action-coverage expansion while penalizing local ambiguity, while the VLA-side score measures informativeness via the frozen Phase-1 policy's no-gradient denoising loss.
CRAPE first filters candidates by buffer-side usefulness and local support, then commits the most VLA-informative remaining candidate, selecting novel yet locally consistent supervision while rejecting redundant or ambiguity-injecting trajectories.

\begin{figure*}[t]
  \centering
  \includegraphics[width=0.95\textwidth]{figure/fig2.png}
    \vspace{-0.15in}
    \caption{
      Overall framework of CRAPE.
      % Phase 1 performs coverage seeding by selecting state-novel subgoals and executing them with a canonical planner, thereby establishing local action references.
      % Phase 2 generates skill-level perturbations around these seeds and acquires candidates that are both buffer-consistent and model-informative.
    }
    \label{fig:introv}
    \vspace{-0.15in}
\end{figure*}

\subsection{State Coverage Seeding via Canonical Subgoal Exploration}
\label{sec:method:m1}
Coverage seeding expands the visited state region while preserving a clean local action reference for later ambiguity assessment. 
CRAPE follows a simple principle in which subgoals are perturbed while the resulting subgoal sequence is executed canonically.
This broadens state support without turning each local region into a multimodal action mixture.

Let $\pi_{\mathrm{code}}$ map the initial state $s_0$ and task context $I$ to a sequence of skill--subgoal pairs:
\begin{equation} \small
\small
\pi_{\mathrm{code}}(s_0, I)
=
\big((z_1,g_1),\ldots,(z_L,g_L)\big).
\label{eq:p1-code-plan}
\end{equation}
For each candidate index $j=1,\ldots,J$, we sample skill-wise subgoal perturbations and form a perturbed subgoal sequence:
\begin{equation} \small
\small
\tilde g_k^{(j)}
=
g_k + \epsilon_{z_k}^{(j)},
\qquad
\epsilon_{z_k}^{(j)} \sim \mathcal D_{z_k},
\qquad
\mathcal G^{(j)}
=
\big((z_1,\tilde g_1^{(j)}),\ldots,(z_L,\tilde g_L^{(j)})\big),
\quad k=1,\ldots,L.
\label{eq:p1-perturbed-subgoals}
\end{equation}
Here, $\mathcal D_{z_k}$ denotes the skill-specific perturbation distribution for skill $z_k$.

To score state coverage, we roll out the canonical planner on each perturbed sequence $\mathcal G^{(j)}$. 
Let $\tilde s_k^{(j)}$ be the state reached after skill $k$, and embed it as
$e_k^{(j)}=\phi_{\mathrm{state}}(z_k,\tilde s_k^{(j)},I)$. 
Its novelty against the corresponding skill-wise seed buffer is computed by a kNN distance:
\begin{equation}
\small
n_s(e,B_I^{(z)})
=
\frac{1}{K_{\mathrm{NN}}}
\sum_{i=1}^{K_{\mathrm{NN}}}
d(e,e_{(i)}),
\label{eq:p1-knn-novelty}
\end{equation}
where $e_{(i)}$ is the $i$-th nearest stored state embedding in $B_I^{(z)}$.
%
We then select the subgoal sequence with the largest aggregate novelty:
\begin{equation} \small
\small
j^\star
=
\operatorname*{arg\,max}_{j\in\{1,\ldots,J\}}
\sum_{k=1}^{L}
n_s(e_k^{(j)},B_I^{(z_k)}),
\qquad
\mathcal G^\star=\mathcal G^{(j^\star)}.
\label{eq:p1-select-subgoal}
\end{equation}
The executable seed trajectory is generated with the canonical policy,
$\xi^\star=\pi_{\mathrm{ca}}(s_0,I,\mathcal G^\star)$. 
By perturbing only subgoals and keeping execution canonical, we expand under-explored state regions while assigning each newly visited region a single low-ambiguity action seed.

\subsection{Consistent Diversity Acquisition}
\label{sec:method:m2}
Given the canonical seeds, we diversify skill executions within the state support, seeking trajectories that differ from prior ones while remaining compatible with local state-conditioned references.

\textbf{Candidate generation.}
For each canonical seed trajectory $\xi_g$ associated with skill $z$ at acquisition step $t$, we generate locally perturbed candidates $\Xi_{I,t}^{(z)}(g) = \{\xi_g^{(1)},\ldots,\xi_g^{(N)}\}$. These perturbations vary the execution of the same skill while preserving feasibility and remaining within the local support established by Phase 1.

\textbf{Action coverage gain.}
Let $B_I^{(z)}$ contain the canonical seeds and accepted acquisition trajectories for skill $z$. For candidate $\xi$, let $A_\xi=(u_1,\ldots,u_T)$ and $\psi_\xi=\phi_{\mathrm{act}}(A_\xi,z)$ be its action-sequence descriptor. We approximate action-coverage gain by descriptor novelty, $\widehat{\Delta H}_A(\mathcal{D}_t,\xi)=n_A(\psi_\xi,B_I^{(z)})$, where
\begin{equation} \small
\small
n_A\!\left(\psi_\xi, B_I^{(z)}\right)
=
\frac{1}{K_{\mathrm{NN}}}
\sum_{i=1}^{K_{\mathrm{NN}}}
d_A\!\left(\psi_\xi, \psi_{(i)}\right).
\label{eq:p2-action-novelty}
\end{equation}
Here, $\psi_{(i)}$ is the $i$-th nearest action-sequence descriptor in $B_I^{(z)}$, and $d_A$ is the descriptor-space distance.

\textbf{Conditional ambiguity penalty.}
Action novelty is useful only when it does not create conflicting actions at already-covered states. For each timestep $\tau$, define the action-chunk descriptor $\psi_{\xi,\tau}=\phi_{\mathrm{act}}(u_{\tau:\tau+h},z)$ and state embedding $e_\tau=\phi_{\mathrm{state}}(z,s_\tau,I)$. Let $\Psi(e_\tau)$ be the action-chunk descriptors of accepted trajectories passing through nearby states under skill $z$, and let $\mathcal{T}_{\mathrm{cov}}(\xi)=\{\tau:|\Psi(e_\tau)|\ge m_{\min}\}$. We estimate the local ambiguity increase as
\begin{equation} \small
\small
\widehat{\Delta H}_{A \mid S}(\mathcal{D}_t,\xi)
=
\frac{1}{|\mathcal{T}_{\mathrm{cov}}(\xi)|}
\sum_{\tau \in \mathcal{T}_{\mathrm{cov}}(\xi)}
\frac{1}{|\Psi(e_\tau)|}
\sum_{\psi \in \Psi(e_\tau)}
d_A\!\left(\psi_{\xi,\tau},\psi\right).
\label{eq:p2-cond-amb}
\end{equation}
A larger value indicates that the candidate deviates from the local action references observed near the same states.

\textbf{Consistent diversity score.}
We combine global action novelty and local consistency as
\begin{equation} \small
M_{\mathrm{MI}}(\mathcal{D}_t,\xi)
=
\beta\,
\widehat{\Delta H}_A(\mathcal{D}_t,\xi)
\;-\;
\lambda\,
\widehat{\Delta H}_{A \mid S}(\mathcal{D}_t,\xi),
\label{eq:p2-mi-score}
\end{equation}
where $\beta$ and $\lambda$ trade off action-support expansion against local ambiguity.

\textbf{Model-side informativeness.}
The buffer-side score tests whether a candidate is distributionally useful, but not whether the current VLA would benefit from observing it. We therefore freeze the Phase-1 policy $\pi_{\theta_1}$ and compute a no-gradient single-step denoising loss on the candidate action chunks, averaged over stochastic denoising evaluations:
\begin{equation} \small
U_{\mathrm{VLA}}(\xi)
\;=\;
\frac{1}{R\,|\mathcal{T}|}
\sum_{r=1}^{R}
\sum_{\tau \in \mathcal{T}}
\ell_{\mathrm{denoise}}\!\left(
\pi_{\theta_1};\; o_\tau,\, I,\, u_{\tau:\tau+h},\, \eta^{(r)}
\right),
\label{eq:p2-model-info}
\end{equation}
where $\eta^{(r)}$ is the $r$-th noise/time sample. A high value means the frozen Phase-1 policy does not yet explain the candidate well, making it informative if it also passes the buffer-side consistency checks.

\textbf{Acquisition rule.}
CRAPE acquires candidates that are both buffer-useful and model-informative. It first forms the feasible set
\begin{equation} \small
\mathcal{F}_{I,t}^{(z)}(g)
=
\left\{
\xi \in \Xi_{I,t}^{(z)}(g)
\;:\;
M_{\mathrm{MI}}(\mathcal{D}_t,\xi) \ge \tau_{\mathrm{MI}},\;
\widehat{\Delta H}_{A \mid S}(\mathcal{D}_t,\xi) \le \tau_{\mathrm{amb}},\;
|\mathcal{T}_{\mathrm{cov}}(\xi)| \ge m_{\min}
\right\},
\label{eq:p2-acquisition}
\end{equation}
then selects $\xi^\star=\operatorname*{arg\,max}_{\xi\in\mathcal{F}_{I,t}^{(z)}(g)}U_{\mathrm{VLA}}(\xi)$. If $\mathcal{F}_{I,t}^{(z)}(g)$ is empty, no trajectory is acquired for this seed. In particular, insufficient local coverage treats the candidate as outside the seeded support and defers acquisition until Phase 1 installs a canonical reference.
%
Let $P_{\mathrm{seed},I}^{(z)}$ be the Phase-1 seed set for skill $z$. The accepted trajectory updates the skill-wise acquisition set and buffer:
\begin{equation} \small
D_{\mathrm{acq},I,t+1}^{(z)}
=
D_{\mathrm{acq},I,t}^{(z)}
\cup
\{\xi^\star\},
\qquad
B_{I,t+1}^{(z)}
=
P_{\mathrm{seed},I}^{(z)}
\cup
D_{\mathrm{acq},I,t+1}^{(z)}.
\label{eq:p2-commit}
\end{equation}
Thus, CRAPE acquires useful diversity: trajectories that expand skill-level action support, remain locally consistent with seeded state-conditioned references, and are informative for the current VLA.

%
%     % 3-1-1: Robo-Scrape소개
%     % 3-1-2: Robo-Scrape component 3개 소개
%     However, such a narrow distribution of unimodal expert behavior may leave out-of-distribution and failure states underrepresented. Thus, scaling purely expert data may not result in a performance gain, but instead risks the generalist overfitting on these data and harming both robustness and generalization
%

% \subsection{Infrastructure for Full Automations}
%     However, such a narrow distribution of unimodal expert behavior may leave out-of-distribution and failure states underrepresented. Thus, scaling purely expert data may not result in a performance gain, but instead risks the generalist overfitting on these data and harming both robustness and generalization\par
%         \textbf{Task-oriented perception module.} ddddd \par
%         \textbf{Phase scheduler with Layout Randomization.} ddddd \par
%         \textbf{Vlm-Based Rollout Validator.} ddddd \par
        
%     \subsection{Hierarchical Perturbation for Improving Task-Relevant Coverage}
% \begin{figure}
%             \centering
%             \includegraphics[width=0.5\linewidth]{placeholder_increase_epi_than_increase_SR.png}
%             \caption{Enter Caption}
%             \label{fig:placeholder}
%         \end{figure}
%                 However, such a narrow distribution of unimodal expert behavior may leave out-of-distribution and failure states underrepresented. Thus, scaling purely expert data may not result in a performance gain, but instead risks the generalist overfitting on these data and harming both robustness and generalization\par
        
%         \textbf{2-layer hierarchy perturbation.}  1. skill-level perturbation, 2. subgoal-level perturbation\par

%         \textbf{Ensemble of Planner as skill-level perturbation.}  \par

%         \textbf{3D gaussian blob as target for sub-goal-level perturbation.}  \par

%     \subsection{Pre-Selective Data Generation for Efficient Scaling}
%         However, such a narrow distribution of unimodal expert behavior may leave out-of-distribution and failure states underrepresented. Thus, scaling purely expert data may not result in a performance gain, but instead risks the generalist overfitting on these data and harming both robustness and generalization
        
%%%%%%%%%%%%%%%%%%%%%%%%%%%%%%%%%%%%%%%%%%%%%%%%%%%%%%%%%%%%%%%%%%%%%%%%%%%%%%%%

%%%%%%%%%%%%%%%%%%%%%%%% EXPERIMENTS %%%%%%%%%%%%%%%%%%%%%%%%%%%%%%%%%%%%%%%%
% All experiment-related floats (Tables 1-6 and figure) live inside \section{Experiments} below.
%===============================================================================

\section{Experiments}
\label{sec:result}
We evaluate whether CRAPE can turn code-generated trajectories into effective real-world supervision for VLA fine-tuning.
The experiments test two central claims:
CRAPE can approach human teleoperation performance and outperform code-based collection baselines.
We further ablate its main components to isolate the roles of state-coverage seeding and pre-selective acquisition.

% \section{Experiments}
% \label{sec:result}
% We evaluate whether CRAPE can turn code-generated trajectories into effective real-world VLA supervision. 
% Our experiments ask four questions: 
% \textbf{(RQ1)} Can CRAPE-collected data train policies that match human teleoperation? 
% \textbf{(RQ2)} Does pre-selective acquisition produce more learnable data than vanilla code-based collection? 
% \textbf{(RQ3)} Which components of CRAPE are necessary? 
% \textbf{(RQ4)} Does the benefit transfer across VLA backbones and robot embodiments?
% RQ3의 OOD/robustness 측면은 Main Results의 jitter 환경 SR과 Ablation의 perturbation 제거 조건에 흡수됨.

%---------------------------------------------------------------
\subsection{Experimental Setup}
\label{sec:setup}

\textbf{Benchmark.}
We construct a real-world benchmark of 10 manipulation tasks covering pick-and-place, distribution, sorting, stacking, articulated lid/drawer manipulation, multi-step arrangement, and deformable bimanual manipulation: 
\textit{pick\_place}~(T1), \textit{distribute\_bread}~(T2), \textit{sort3}~(T3), \textit{stack3}~(T4), \textit{close\_pot}~(T5), \textit{open\_drawer}~(T6), \textit{prepare\_meal}~(T7), \textit{fold\_towel}~(T8), \textit{open\_lid}~(T9), and \textit{close\_drawer}~(T10). 
The main benchmark uses the SO-101 platform, with a dual-SO-101 setup.
% ; cross-embodiment results on UR7e are reported in Table~\ref{tab:crossembod}.

\textbf{Baselines.}
We compare CRAPE against three data collection baselines: human teleoperation; RoboTwin~\citep{chen2025robotwin}, a code-based simulation pipeline whose code-generation scheme we port to real-world collection; and Manipulate-Anything (MA)~\citep{duan2024manipulate}, a VLM-driven open-loop baseline.

\textbf{Training and evaluation.}
All methods use the same LLM-generated primitive skill library and the same Isaac~Lab~\citep{mittal2025isaac} sim-pretraining stage before real-world fine-tuning. 
We evaluate two VLA backbones, smolVLA~\citep{shukor2025smolvla} and GR00T~N1.5~\citep{bjorck2025gr00t,nvidiagear2025gr00tn15}, each fine-tuned per task under the same protocol. 
Each trained policy is evaluated over 30 rollouts under matched initial conditions using task-specific success criteria. 
Fine-tuning consumed approximately \textit{250} GPU-hours on H100 GPUs.

% =================================================================
% [Table 1] Main Results: 4 collection methods x 2 VLAs x 10 tasks
% =================================================================
\begin{table*}[t]
\centering
\caption{Main results on the real-world benchmark. We measure success rate (\%) over 10 manipulation tasks for two VLA backbones, trained under the same fine-tuning procedure across ours and the baselines.}
\label{tab:main}
\begin{adjustbox}{width=0.98\textwidth}
\small
\setlength{\tabcolsep}{3pt}
\begin{tabular}{l|l | cccccccccc | c}
\toprule
\textbf{Method} & \textbf{VLA}
& \textbf{T1} & \textbf{T2} & \textbf{T3} & \textbf{T4} & \textbf{T5}
& \textbf{T6} & \textbf{T7} & \textbf{T8} & \textbf{T9} & \textbf{T10}
& \textbf{Overall} \\
\midrule
\multirow{2}{*}{Teleoperation}
  & smolVLA     & 83.33 & 66.67 & 0 & 0 & 6.67 & 96.67 & 50.00 & 93.33 & 93.33 & 93.33 & 58.33 \\
  & GR00T N1.5  & 40.00 & 26.67 & 0 & 0 & 0 & 76.67 & 36.67 & 23.33 & 0 & 100 & 30.33 \\
\midrule
\multirow{2}{*}{RoboTwin}
  & smolVLA     & 0 & 0 & 0 & 0 & 56.67 & 100 & 3.33 & 0 & 6.67 & 0 & 16.67 \\
  & GR00T N1.5  & 13.33 & 0 & 0 & 0 & 0 & 96.67 & 23.33 & 0 & 0 & 13.33 & 14.67 \\
\midrule
\multirow{2}{*}{MA}
  & smolVLA     & 0 & 0 & 13.33 & 0 & 63.33 & 100 & 0 & 0 & 0 & 0 & 17.67 \\
  & GR00T N1.5  & 6.67 & 0 & 0 & 0 & 0 & 93.33 & 30.00 & 0 & 0 & 26.67 & 15.67 \\
\midrule
\multirow{2}{*}{\textbf{CRAPE (Ours)}}
  & smolVLA     & 100 & 40 & 33.33 & 6.67 & 33.33 & 100 & 96.67 & 83.33 & 100 & 86.67 & 68.00 \\
  & GR00T N1.5  & 23.33 & 6.67 & 0 & 0 & 53.33 & 100 & 100 & 46.67 & 0 & 93.33 & 42.33 \\
\bottomrule
\end{tabular}
\end{adjustbox}
\end{table*}

%---------------------------------------------------------------
% \subsection{Main Results}
% \label{sec:exp_main}

% The main results shows whether CRAPE-collected data can serve as effective real-world supervision for VLA fine-tuning, matching human teleoperation while outperforming vanilla code-based collection. 
% Table~\ref{tab:main} reports success rates across 10 tasks, four collection methods, and two VLA backbones under the same sim-pretrain and real-world fine-tune protocol. 
% All methods use an identical Isaac~Lab sim-pretraining stage before real-world fine-tuning, which equalizes the starting point across methods and isolates the effect of real-world data acquisition.

% \textbf{Teleoperation-level performance.}
% CRAPE achieves success rates comparable to policies trained on human teleoperation across the benchmark. 
% This indicates that code-generated trajectories can replace much of the manual demonstration burden when they are selectively acquired as clean and informative VLA supervision.

% \textbf{Comparison to vanilla code-based collection.}
% RoboTwin~2.0 and Manipulate-Anything often generate feasible real-world rollouts, but their unfiltered trajectories are less effective for closed-loop VLA training. 
% Policies trained on these data are more brittle under actuation jitter and small state deviations, suggesting that executable demonstrations are not necessarily learnable demonstrations. 
% CRAPE improves over these baselines by selecting trajectories that expand coverage while remaining locally consistent with the existing buffer.

% \textbf{Consistency across VLA backbones.}
% The same trend holds for both smolVLA and GR00T~N1.5, showing that CRAPE's gains are not tied to a specific policy architecture. 
% Rather, the improvement comes from the quality of the acquired data: state-coverage seeding provides clean support, and pre-selective acquisition admits candidates that are diverse, consistent, and informative.
% % \subsection{Main Results}
% % \label{sec:exp_main}

% % Table~\ref{tab:main} reports per-task and overall SR for all four collection
% % methods across the two VLA backbones, targeting three aims:
% % (i)~SR competitiveness against human teleoperation,
% % (ii)~superiority over vanilla code-based collection (RoboTwin, MA), and
% % (iii)~generality across off-the-shelf VLA backbones.

% % \textbf{Protocol specific to this experiment.}
% % All methods share an identical sim-pretraining stage on Isaac~Lab
% % demonstrations before real-world fine-tuning; without it, RoboTwin- and
% % MA-collected data alone drive the policy into states from which it cannot
% % recover, so sim-pretraining equalizes the floor across methods.

% % \textbf{(RQ1) CRAPE matches human teleoperation.}
% % Across the 10 tasks, policies trained on CRAPE-collected data reach overall
% % SR comparable to human teleoperation.

% % \textbf{(RQ2) Vanilla code-based collection fails to transfer to the real world.}
% % Despite the shared sim-pretraining stage, RoboTwin- and MA-trained policies yield
% % low SR (overall \textit{XX\%}); SO-101's actuation jitter repeatedly drives
% % the policy out-of-distribution. Vanilla code-based collection is therefore
% % unfit for real-world VLA training.

% % \textbf{(RQ3) The CRAPE recipe generalizes across VLAs.}
% % Phase~1 canonical subgoal seeding broadens state coverage, and Phase~2
% % pre-selective acquisition admits only useful trajectories, producing a VLA
% % that is robust to jitter-induced OOD while remaining consistent in its
% % predictions. The pattern is preserved across both smolVLA and GR00T~N1.5,
% % confirming that the advantage is a property of the collected data.

% % =================================================================
% %========================== CHANGED ↓ : Data Quality (scatter + trend figures + narrative; replaces former Table 2) ==========================
% % Data Quality & Sample Efficiency — visualized as two figures:
% %   fig:quality_scatter (100ep snapshot, diversity vs NCD)
% %   fig:quality_trend   (scaling over 20-100 episodes, mean +/- std over tasks)
% % =================================================================

% \begin{figure}[t]
%   \centering
%   \includegraphics[width=\linewidth]{figure/traj_visualization.png}
%   \caption{Per-task diversity (Signature) vs.\ diversity-invariant consistency (NCD) at $100$ episodes. Each marker is a task; shading is per-method hdensity. Higher on the plot is more consistent.}
%   \label{fig:traj_visualization}
%   \vspace{-0.1in}
% \end{figure}



% \begin{figure}[t]
%   \centering
%   \includegraphics[width=0.66\linewidth]{figure/fig_quality_scatter.png}
%   \caption{Per-task diversity (Signature) vs.\ diversity-invariant consistency (NCD) at $100$ episodes. Each marker is a task; shading is per-method hdensity. Higher on the plot is more consistent.}
%   \label{fig:quality_scatter}
%   \vspace{-0.1in}
% \end{figure}

% \begin{figure*}[t]
%   \centering
%   \includegraphics[width=\textwidth]{figure/fig_quality_trend.png}
%   \caption{Dataset quality vs.\ acquisition budget. Lines: mean over the $8$ single-arm tasks; bands: $\pm 1$ std across tasks.}
%   \label{fig:quality_trend}
%   \vspace{-0.15in}
% \end{figure*}

% %---------------------------------------------------------------
% \subsection{Data Quality and Sample Efficiency}
% \label{sec:exp_quality}

% Figures~\ref{fig:quality_scatter} and~\ref{fig:quality_trend} ask whether CRAPE's gains come from collecting \emph{more} episodes or \emph{more learnable} ones, tracking diversity, coverage, and NCD as the acquisition budget grows ($20\!\to\!100$ episodes).

% CRAPE is the only method whose dataset quality and SR improve together with scale. It raises Signature Diversity and Transition Coverage toward teleoperation level (Ours $10.6$/$0.5$ vs.\ RoboTwin $7.3$/$\!-\!7.5$ at $100$ episodes) while keeping conditional ambiguity low \emph{relative to its diversity}. RoboTwin saturates under the same budget: its deterministic primitives make extra rollouts redundant, so executable rollouts are not learnable supervision.

% Because diversity and raw consistency trade off, we measure consistency with NCD (local ambiguity normalized by global diversity). Teleoperation is diverse but comparatively inconsistent---heterogeneous human action modes give it high conditional ambiguity---whereas CRAPE attains \emph{lower} ambiguity on all $8$ single-arm tasks (NCD; Wilcoxon $p=0.004$, mean gap $0.063$, $95\%$ CI $[0.034,0.091]$). Figure~\ref{fig:quality_scatter} summarizes the trade-off at $100$ episodes: RoboTwin occupies the low-diversity region, teleoperation the high-diversity/high-ambiguity region, and CRAPE the desired high-diversity/low-ambiguity ``useful-diversity'' regime. These advantages strengthen with the acquisition budget (Figure~\ref{fig:quality_trend}): CRAPE's diversity keeps growing while RoboTwin saturates, and NCD improves at every budget with CRAPE more consistent than teleoperation throughout---CRAPE collects \emph{more learnable}, not merely more, data.

\subsection{Main Results}
\label{sec:exp_main}

The main results show whether CRAPE-collected data can serve as effective real-world supervision for VLA fine-tuning, matching human teleoperation while outperforming vanilla code-based collection.
Table~\ref{tab:main} reports success rates across 10 tasks, four collection methods, and two VLA backbones under the same simulation pre-training and real-world fine-tuning protocol.
All methods use an identical Isaac~Lab sim-pretraining stage before real-world fine-tuning, which equalizes the starting point across methods and isolates the effect of real-world data acquisition.

\textbf{Teleoperation-level performance.}
CRAPE matches and, in macro-average, exceeds human teleoperation on both backbones ($68.0$ vs.\ $58.33$ on smolVLA, $42.33$ vs.\ $30.33$ on GR00T~N1.5).
This indicates that selectively acquired code-generated trajectories can replace much of the manual demonstration burden and serve as effective supervision in the real world.

\textbf{Comparison to vanilla code-based collection.}
RoboTwin~2.0 and Manipulate-Anything generate feasible rollouts, but their unfiltered trajectories are far less learnable: CRAPE outperforms the stronger of the two by $\sim\!3.9\times$ on smolVLA ($68.0$ vs.\ $17.67$) and $\sim\!2.7\times$ on GR00T ($42.33$ vs.\ $15.67$).
Their deterministic primitives accumulate stereotyped trajectories, so success concentrates on primitive-aligned tasks like \textit{open\_drawer}~(T6; $\ge\!93\%$) while the policies cover little state space and drift out of distribution on free-space tasks like \textit{pick\_place}~(T1; $\le\!13.33$, where CRAPE reaches $100$).
By selecting trajectories that expand coverage while staying locally consistent, CRAPE confirms that executable demonstrations are not necessarily learnable ones.

\textbf{Consistency across VLA backbones.}
The ordering CRAPE~$>$~teleoperation~$>$~{RoboTwin,~MA} holds on both smolVLA and GR00T~N1.5, showing that the gains are not architecture-specific.
They instead come from better data: state-coverage seeding builds clean support, and pre-selective acquisition admits diverse, consistent, and informative candidates.
% \subsection{Main Results}
% \label{sec:exp_main}

% Table~\ref{tab:main} reports per-task and overall SR for all four collection
% methods across the two VLA backbones, targeting three aims:
% (i)~SR competitiveness against human teleoperation,
% (ii)~superiority over vanilla code-based collection (RoboTwin, MA), and
% (iii)~generality across off-the-shelf VLA backbones.

% \textbf{Protocol specific to this experiment.}
% All methods share an identical sim-pretraining stage on Isaac~Lab
% demonstrations before real-world fine-tuning; without it, RoboTwin- and
% MA-collected data alone drive the policy into states from which it cannot
% recover, so sim-pretraining equalizes the floor across methods.

% \textbf{(RQ1) CRAPE matches human teleoperation.}
% Across the 10 tasks, policies trained on CRAPE-collected data reach overall
% SR comparable to human teleoperation.

% \textbf{(RQ2) Vanilla code-based collection fails to transfer to the real world.}
% Despite the shared sim-pretraining stage, RoboTwin- and MA-trained policies yield
% low SR (overall \textit{XX%}); SO-101's actuation jitter repeatedly drives
% the policy out-of-distribution. Vanilla code-based collection is therefore
% unfit for real-world VLA training.

% \textbf{(RQ3) The CRAPE recipe generalizes across VLAs.}
% Phase~1 canonical subgoal seeding broadens state coverage, and Phase~2
% pre-selective acquisition admits only useful trajectories, producing a VLA
% that is robust to jitter-induced OOD while remaining consistent in its
% predictions. The pattern is preserved across both smolVLA and GR00T~N1.5,
% confirming that the advantage is a property of the collected data.

% =================================================================
%========================== CHANGED ↓ : Data Quality (scatter + trend figures + narrative; replaces former Table 2) ==========================
% Data Quality & Sample Efficiency — visualized as two figures:
%   fig:quality_scatter (100ep snapshot, diversity vs NCD)
%   fig:quality_trend   (scaling over 20-100 episodes, mean +/- std over tasks)
% =================================================================

\begin{figure}[t]
\centering
\includegraphics[width=\linewidth]{figure/traj_visualization.png}
\caption{
Dataset quality at $100$ episodes. Left: per-task Signature Diversity and normalized conditional divergence (NCD); the lower-right region indicates high diversity with low conditional ambiguity. Shading denotes per-method density. Right: representative end-effector trajectories and rollout frames.
}
\label{fig:quality_scatter}
\vspace{-0.1in}
\end{figure}

\begin{figure*}[t]
\centering
\includegraphics[width=\textwidth]{figure/fig_quality_trend.png}
\caption{
Dataset quality as the acquisition budget increases. Lines show the mean over the $8$ single-arm tasks, and shaded bands show $\pm 1$ standard deviation. Transition Coverage is a standardized support score, so negative values indicate below-reference coverage.
}
\label{fig:quality_trend}
\vspace{-0.1in}
\end{figure*}

%---------------------------------------------------------------
\subsection{Data Quality}
\label{sec:exp_quality}

We ask whether more acquisition yields useful variation or merely more rollouts. At $100$ episodes (Figure~\ref{fig:quality_scatter}), RoboTwin is low-diversity and teleoperation is high-diversity but high-ambiguity, whereas CRAPE occupies the desired high-diversity, low-ambiguity region: it beats RoboTwin on Signature Diversity and Transition Coverage ($10.6$/$0.5$ vs.\ $7.3$/$-7.5$) and attains lower NCD than teleoperation on all $8$ single-arm tasks\footnote{T8 (bimanual) and T10 (corrupted logs) are excluded; the eight are T1--T7 and T9.} (Wilcoxon $p=0.004$). CRAPE thus collects more learnable data, not just more, consistent with its teleoperation-level success in Table~\ref{tab:main}.

With more budget (Figure~\ref{fig:quality_trend}), CRAPE keeps improving: Signature Diversity rises, Transition Coverage stays far above RoboTwin, and NCD falls below teleoperation. RoboTwin instead saturates, its deterministic primitives making its low NCD a trivial artifact of low diversity rather than learnable supervision.



% \subsection{Data Quality and Sample Efficiency}
% \label{sec:exp_quality}

% Table~\ref{tab:quality} tracks four dataset-level quality metrics and downstream SR as episodes grow ($30/50/100$) over a common pool of 10 fixed layouts. The central claim is \emph{joint}: useful diversity means expanding action coverage while keeping local actions consistent, which only CRAPE achieves.

% \textbf{Metrics.}
% \textit{Signature Diversity}~\citep{faktual2026} measures episode-level trajectory variation.
% \textit{Transition Support Coverage} measures how widely the data covers local $[s,\Delta a,s']$ transitions.
% \textit{Local Action Inconsistency}~\citep{belkhale2023data} measures conditional action ambiguity at similar states (lower is better; legitimate multimodality is not penalized).
% \textit{Mutual Information}~\citep{hejna2025robot} between state and future-action descriptors is an auxiliary summary, reported because it coincides with CRAPE's selection objective.

% \textbf{Diversity vs.\ consistency.}
% CRAPE raises Signature Diversity and Transition Coverage to Teleoperation level while RoboTwin/MA stay flat (redundant rollouts add little coverage). Unlike Teleoperation, whose diversity comes with high Action Inconsistency, CRAPE keeps inconsistency low, expanding coverage without injecting the action ambiguity that destabilizes flow-/diffusion-based VLAs.

% \textbf{Scaling.}
% CRAPE's quality metrics and SR improve together with episodes, whereas RoboTwin/MA plateau at low SR despite matching dataset size; Mutual Information co-moves with SR, indicating the dataset-level signal predicts learnability. The diversity--inconsistency view (Figure~\ref{fig:quality_scatter}) makes this explicit: Teleoperation occupies the high-diversity/high-inconsistency corner, RoboTwin/MA the low/low corner, and CRAPE alone the high-diversity/low-inconsistency ``useful'' corner.
% % TODO: Figure fig:quality_scatter (x=Signature Diversity, y=Local Action Inconsistency, 30->50->100 화살표). docs/table2_experiments.md §6.2 참고.

% %---------------------------------------------------------------
% \subsection{Data Quality and Sample Efficiency}
% \label{sec:exp_quality}

% Table~\ref{tab:quality} tracks five dataset-level metrics as the number of collected episodes grows (30, 50, 100). Figure~\ref{fig:placeholder_increase_epi_than_increase_SR} visualizes the corresponding SR-vs-episode scaling curves.

% \textbf{Metric definitions.}
% \textit{Diversity} is measured as pairwise trajectory distance (DTW) averaged across collected episodes.
% \textit{Action consistency} is the per-state action variance, with multi-modal modes treated as valid (not penalized) following diffusion/flow-based VLA assumptions.
% \textit{State coverage} is the volume of the visited state distribution (kernel-density estimate).
% \textit{Info.\ gain} is the policy-uncertainty reduction proxy used by our pre-selective acquisition module.
% % 위 정의는 method 또는 supplementary로 더 정밀하게 빼도 됨.

% \textbf{(RQ4-i) Sample efficiency.}
% % 동일 SR 도달까지 필요한 episode 수에서 CRAPE의 우위, 또는 동일 episode 수에서의 SR 격차 서술.

% \textbf{(RQ4-ii) Pre-selective vs.\ post-hoc filtering.}
% % Pre-selective ON / OFF / Post-hoc filtering 비교가 ablation(Table 3)과 본 표 결합으로 어떻게 입증되는지 서술.
% % 핵심 메시지: post-hoc은 robot time을 환불받지 못한다.

% =================================================================
% [Table 3] Ablation Study
% =================================================================
\begin{table*}[t]
\centering
\caption{Cumulative leave-out ablation on \textit{pick\_place}: from \textbf{Ours} we remove one component at a time down to the vanilla CaP baseline.}
\small
\setlength{\tabcolsep}{6pt}
\begin{tabular}{l cccc ccc}
\toprule
& \multicolumn{4}{c}{\textbf{Component}} & \multicolumn{3}{c}{\textbf{Metric}} \\
\cmidrule(lr){2-5}\cmidrule(lr){6-8}
\textbf{Variant}
& Subgoal Seed.\ & Traj.\ Perturb.\ & $M_{\mathrm{MI}}$ & $U_{\mathrm{VLA}}$
& \textbf{SR}\,$\uparrow$ & SigDiv\,$\uparrow$ & TransCov\,$\uparrow$ \\
\midrule
\textbf{Ours}            & \checkmark & \checkmark & \checkmark & \checkmark & \textbf{96.7} & 11.0 & 2.7 \\
+w/o $U_{\mathrm{VLA}}$   & \checkmark & \checkmark & \checkmark & --         & 93.3 & 10.6 & 2.5 \\
+w/o $M_{\mathrm{MI}}$    & \checkmark & \checkmark & --         & --         & 80.0 & 10.8 & 3.1 \\
+w/o Traj.\ Perturb.      & \checkmark & --         & --         & --         & 90.0 & 10.0 & 3.2 \\
+w/o Subgoal Seed. & --         & --         & --         & --         & 13.3 & 9.3  & $-4.6$ \\
\bottomrule
\end{tabular}
\label{tab:ablation}
\end{table*}

\subsection{Ablation Studies}
\label{sec:exp_ablation}

We ablate CRAPE on \textit{pick\_place} by removing one component at a time from \textbf{Ours} down to the vanilla CaP baseline (Table~\ref{tab:ablation}).
Removing subgoal seeding is the most damaging: SR falls from $90.0\%$ to $13.3\%$ and Transition Coverage from $3.2$ to $-4.6$, as the policy is no longer seeded across the states it must reach at test time.
Removing the buffer-side score (w/o $M_{\mathrm{MI}}$) then admits the perturbed candidates at random and lowers SR to $80.0\%$, below the seed-only variant (w/o Traj.\ Perturb., $90.0\%$) despite its higher Signature Diversity ($10.0\!\to\!10.8$), because indiscriminate admission injects locally inconsistent, ambiguity-increasing demonstrations.
Reinstating $M_{\mathrm{MI}}$ filters these and recovers SR to $93.3\%$, and adding the model-side score $U_{\mathrm{VLA}}$ (\textbf{Ours}) reaches $96.7\%$.

%---------------------------------------------------------

% =================================================================
% [Table 4] Cross-embodiment Generalization
% =================================================================
% \begin{table}[t]
% \centering
% \caption{Cross-embodiment generalization. Success rate (\%) on 3 tasks for an alternative embodiment, comparing RoboTwin-collected and CRAPE-collected data.}
% \label{tab:crossembod}
% \small
% \begin{tabular}{llcccc}
% \toprule
% \textbf{Method} & \textbf{Task1} & \textbf{Task2} & \textbf{Task3} & \textbf{Overall} \\
% \midrule
% RoboTwin & 20.00 & 6.67 & 0.00 & 8.89 \\
% \textbf{CRAPE (Ours)} & 46.67 & 20.00 & 6.67 & 24.45 \\
% \bottomrule
% \end{tabular}
% \end{table}

% %---------------------------------------------------------------
% \subsection{Cross-Embodiment Generalization}
% \label{sec:exp_crossembod}

% Table~\ref{tab:crossembod} tests whether CRAPE transfers beyond the SO-101. Porting the identical pipeline to a UR7e arm, CRAPE beats RoboTwin on every task and nearly triples overall SR ($8.89\%\!\to\!24.45\%$). Given the same skill library, the framework applies unchanged to this different embodiment and yields the same advantage, indicating that the gain comes from \emph{how} trajectories are selected (useful-diversity-guided coverage expansion) and not from a specific robot.

% 결과 채워지면, 본 절에서 RoboTwin 대비 격차 및 SO-101 대비 절대 SR 변화 서술.

%---------------------------------------------------------------
% Supplementary로 빼는 항목 (본문에서 1줄 참조):
% Pipeline-level cost (episodes/hour, autonomous duty cycle, human interventions/episode) — Supplementary Table~X.

%===============================================================================


\section{Conclusion}
\label{sec:conclusion}

We presented CRAPE, a code-based real-world data-acquisition framework that treats automated collection as \emph{pre-selective acquisition} for VLA fine-tuning. 
Rather than adding all generated rollouts, CRAPE first builds clean state--action support through canonical subgoal exploration and then admits candidates that expand action coverage, remain locally consistent, and are informative to the current VLA. 
Across 10 SO-101 tasks, CRAPE reaches teleoperation-level performance, outperforms vanilla code-based collection under the same protocol, and ablations show that random admission increases ambiguity and reduces success rates.

\textbf{Limitation.} CRAPE still relies on a pre-defined skill library and a canonical planner, and extending it to less structured, longer-horizon, open-ended manipulation remains future work.%===============================================================================

% =================================================================================
% \section{Limitations}
% All submissions should include a Limitations section (counted toward the 8-page limit), explicitly describing limiting assumptions, failure modes, and other limitations of the results and experiments, and how these might be addressed in the future


% \section{Citations}
% \label{sec:citations}

% 	Citations can be made using either \textbackslash citep\{\} or \textbackslash citet\{\}, depending from the appropriateness. To avoid the citation moving to the next line, it is often a good practice to replace the space before with a tilde (\~{}) character.
% 	Example 1: ``CoRL is the best conference ever~\citep{Gauss1857}.''
% 	Example 2: ``\citet{Lagrange1788} proved, both theoretically and numerically, that CoRL is the best conference ever.''

\clearpage
% The acknowledgments are automatically included only in the final and preprint versions of the paper.
\acknowledgments{If a paper is accepted, the final camera-ready version will (and probably should) include acknowledgments. All acknowledgments go at the end of the paper, including thanks to reviewers who gave useful comments, to colleagues who contributed to the ideas, and to funding agencies and corporate sponsors that provided financial support.}

%===============================================================================

% no \bibliographystyle is required, since the corl style is automatically used.
\bibliography{example}  % .bib

%===============================================================================
\clearpage
\appendix
\renewcommand{\thetable}{A\arabic{table}}
\renewcommand{\thefigure}{A\arabic{figure}}
\setcounter{table}{0}
\setcounter{figure}{0}


% ==============================================================================
\section{Dataset Quality Metrics}
\label{app:metrics}

We define the three metrics used in Section~\ref{sec:exp_quality} (Figures~\ref{fig:quality_scatter}, \ref{fig:quality_trend}) and Table~\ref{tab:ablation}: Signature Diversity, Normalized Conditional Divergence (NCD), and Transition Coverage. Each episode is the $6$ DoF end effector path $\{\mathbf{s}_t\!\in\!\mathbb{R}^{6}\}_{t=1}^{T}$ from forward kinematics, with $\Delta\mathbf{s}_t\!=\!\mathbf{s}_{t+1}\!-\!\mathbf{s}_t$. All values are per task on the committed buffer and aggregated over the eight single arm tasks.

\textbf{Signature Diversity (SigDiv, $\uparrow$).}
Let $\mathbf{K}_{\mathrm{sig}}\!\in\!\mathbb{R}^{N\times N}$ be the signature kernel Gram matrix over the $N$ episode paths~\citep{faktual2026}, and $\{\lambda_i\}$ the eigenvalues of its trace normalized form ($\sum_i\lambda_i\!=\!1$):
\begin{equation} \small
\mathrm{SigDiv}(\mathcal{D}) \;=\; \exp\!\Bigl(-\textstyle\sum_{i}\lambda_i\log\lambda_i\Bigr).
\label{eq:metric-sigdiv}
\end{equation}
This is the exponentiated spectral entropy of $\mathbf{K}_{\mathrm{sig}}$, i.e., the effective trajectory rank in $[1, N]$.

\textbf{Normalized Conditional Divergence (NCD, $\downarrow$).}
Let $\mathcal{N}_k(\mathbf{s}_t)$ be the $k\!=\!20$ state neighbors of $\mathbf{s}_t$ drawn from other episodes, $d_{\min}(t)\!=\!\min_{u\in\mathcal{N}_k(\mathbf{s}_t)}\lVert\Delta\mathbf{s}_t-\Delta\mathbf{s}_u\rVert$ the mode aware nearest residual, and $\bar d_{\mathrm{glob}}$ the mean action distance over pairs from different episodes in $\mathcal{D}$:
\begin{equation} \small
\mathrm{NCD}(\mathcal{D}) \;=\; \frac{\mathbb{E}_t\bigl[\,d_{\min}(t)\,\bigr]}{\bar d_{\mathrm{glob}}}.
\label{eq:metric-ncd}
\end{equation}
NCD is a scale free ratio. Deterministic baselines collapse $d_{\min}$ toward zero and yield a trivially low NCD that is an artifact of narrow support rather than meaningful local consistency, so NCD is read jointly with SigDiv.

\textbf{Transition Coverage (TransCov, $\uparrow$).}
The fingerprint $\tilde{\mathbf{x}}\!=\!(\mathbf{s},\Delta\mathbf{s})\!\in\!\mathbb{R}^{12}$ is standardized per task per dimension to zero mean and unit variance, then we report its Kozachenko Leonenko differential entropy:
\begin{equation} \small
\mathrm{TransCov}(\mathcal{D}) \;=\; \hat H_{\mathrm{KL}}\!\bigl[\tilde{\mathbf{x}}\bigr]\quad[\mathrm{nats}].
\label{eq:metric-transcov}
\end{equation}
Higher values indicate wider transition support; the score becomes negative when the standardized fingerprint concentrates below the unit volume baseline implicit in differential entropy. This per task standardization makes cross task averaging in Figure~\ref{fig:quality_trend} valid.

\textbf{Aggregation.} Trend figures report per task means with $\pm 1$ std bands across the eight tasks. Uncertainty uses episode subsampling bootstrap ($B\!=\!200$, $70\%$ without replacement). Paired tests use Wilcoxon signed rank; at $100$ episodes Ours yields lower NCD than teleoperation on $8/8$ tasks ($p\!=\!0.004$).

% ==============================================================================
\section{Fairness of the Teleoperation Baseline}
\label{app:teleop}

A skeptic might read teleoperation's low scores in Table~\ref{tab:main} (notably
$0\%$ on \textit{sort3} (T3) and \textit{stack3} (T4)) as an under resourced
baseline rather than competitive code generated data. We rule this out along two
axes: training budget parity (B.1) and extra teleoperation data (B.2).

\textbf{(B.1) Training budget parity.}
Within each backbone, every method uses the same fine tuning recipe
(Table~\ref{tab:budget_parity}); only the trajectory source differs. Epochs are
fixed per backbone, so optimizer iterations match across methods for any task.

\begin{table}[h]
\centering
\caption{Training budget parity. Within each backbone all values are constant across the $10$ tasks and four methods; only the trajectory source differs. Both backbones use AdamW with a cosine schedule with warmup.}
\label{tab:budget_parity}
\small
\setlength{\tabcolsep}{6pt}
\resizebox{\textwidth}{!}{%
\begin{tabular}{llccccc}
\toprule
\textbf{Backbone} & \textbf{Method} & \textbf{Committed ep.} & \textbf{Epochs} & \textbf{LR} & \textbf{Global batch} & \textbf{Sim pretrain} \\
\midrule
\multirow{4}{*}{smolVLA}
  & Teleoperation       & 100 & 50 & $1\times10^{-4}$ $\rightarrow$ $2.5\times10^{-6}$ & 256 & shared \\
  & RoboTwin            & 100 & 50 & $1\times10^{-4}$ $\rightarrow$ $2.5\times10^{-6}$ & 256 & shared \\
  & Manipulate-Anything & 100 & 50 & $1\times10^{-4}$ $\rightarrow$ $2.5\times10^{-6}$ & 256 & shared \\
  & \textbf{CRAPE (Ours)} & 100 & 50 & $1\times10^{-4}$ $\rightarrow$ $2.5\times10^{-6}$ & 256 & shared \\
\midrule
\multirow{4}{*}{GR00T~N1.5}
  & Teleoperation       & 100 & 100 & $1\times10^{-4}$ $\rightarrow$ $1\times10^{-5}$ & 256 & shared \\
  & RoboTwin            & 100 & 100 & $1\times10^{-4}$ $\rightarrow$ $1\times10^{-5}$ & 256 & shared \\
  & Manipulate-Anything & 100 & 100 & $1\times10^{-4}$ $\rightarrow$ $1\times10^{-5}$ & 256 & shared \\
  & \textbf{CRAPE (Ours)} & 100 & 100 & $1\times10^{-4}$ $\rightarrow$ $1\times10^{-5}$ & 256 & shared \\
\bottomrule
\end{tabular}
}
\end{table}


\textbf{(B.2) Teleoperation with extra data.}
We grant teleoperation up to $2\times$ its matched budget on its two $0\%$ tasks,
with all hyperparameters fixed (Table~\ref{tab:teleop_saturation}). The zeros
reflect data scarcity rather than tuning, since teleoperation improves with
budget. Even so, on \textit{sort3} CRAPE at $100$ episodes ($33.33\%$) exceeds
teleoperation at $200$ episodes ($26.67\%$), reaching higher success at half the
data; \textit{stack3} is a CRAPE limitation (Section~\ref{sec:conclusion}). We
read this table as a data efficiency check, not a matched comparison.

\begin{table}[h]
\centering
\caption{Teleoperation success rate (\%) on its zero success tasks as the budget scales to $2\times$ (smolVLA), hyperparameters fixed. The CRAPE column is at its own $100$ episode budget, not a matched comparison.}
\label{tab:teleop_saturation}
\small
\setlength{\tabcolsep}{6pt}
\begin{tabular}{lccc|c}
\toprule
& \multicolumn{3}{c|}{\textbf{Teleoperation budget}} & \textbf{CRAPE} \\
\cmidrule(lr){2-4}
\textbf{Task} & \textbf{100 ep (Table~\ref{tab:main})} & \textbf{150 ep} & \textbf{200 ep} & \textbf{(ref., 100 ep)} \\
\midrule
\textit{sort3} (T3)  & 0 & 6.67 & 26.67 & 33.33 \\
\textit{stack3} (T4) & 0 & 13.33 & 20.00 & 6.67 \\
\bottomrule
\end{tabular}
\end{table}

% ==============================================================================
\section{Effect of Simulation-Stage Pretraining}
\label{app:sim2real}

All methods in Table~\ref{tab:main} are sim-pretrained on Isaac~Lab before
real-world fine-tuning. Table~\ref{tab:sim2real} ablates this stage by training
the same real-world data either directly (\emph{Only-real}) or after sim
pretraining (\emph{Sim2real}). Without the sim stage, code-based collection
(RoboTwin) collapses to $\sim\!0\%$: its rollouts drive the policy into
unrecoverable out-of-distribution states from which no task is completed, whereas
human teleoperation---already in-distribution---degrades far less. The sim stage
installs a basic in-distribution prior that makes code-based real-world data
trainable, which is why it is shared across all methods in the main results.

\begin{table}[h]
\centering
\caption{Success rate (\%) with and without the simulation pretraining stage (smolVLA). \emph{Only-real} fine-tunes the base checkpoint directly; \emph{Sim2real} is the main-paper protocol. $0^{\dagger}$ denotes rollouts that terminate in unrecoverable OOD states.}
\label{tab:sim2real}
\small
\setlength{\tabcolsep}{6pt}
\begin{tabular}{llcccc}
\toprule
\textbf{Method} & \textbf{Sim stage} & \textit{pick\_place} & \textit{distribute\_bread} & \textit{close\_pot} & \textbf{Overall} \\
\midrule
\multirow{2}{*}{Teleoperation}
  & Only-real & 53.33 & $0^{\dagger}$ & $0^{\dagger}$ & 17.78 \\
  & Sim2real  & 83.33 & 66.67 & 6.67 & 52.22 \\
\midrule
\multirow{2}{*}{RoboTwin}
  & Only-real & $0^{\dagger}$ & $0^{\dagger}$ & $0^{\dagger}$ & $0^{\dagger}$ \\
  & Sim2real  & $0^{\dagger}$ & $0^{\dagger}$ & 56.67 & 18.89 \\
\midrule
\multirow{2}{*}{\textbf{CRAPE (Ours)}}
  & Only-real & 80.00 & 6.67 & $0^{\dagger}$ & 28.89 \\
  & Sim2real  & 100 & 40.00 & 33.33 & 57.78 \\
\bottomrule
\end{tabular}
\end{table}

% ==============================================================================
\section{Cross-Embodiment Generalization}
\label{app:crossembod}

To test whether CRAPE's gain is specific to the SO-101, we port the same
pipeline to a UR7e arm and compare code-based collection (RoboTwin) against CRAPE
on \textit{pick\_place}, \textit{distribute\_bread}, and \textit{sort3}
(Table~\ref{tab:crossembod}). CRAPE beats CaP on every task and nearly triples
overall SR ($8.89\%\!\to\!24.45\%$), indicating that the advantage comes from how
trajectories are selected rather than from the embodiment.

\begin{table}[h]
\centering
\caption{Cross-embodiment generalization. Success rate (\%) on an alternative embodiment (UR7e), comparing RoboTwin-collected and CRAPE-collected data.}
\label{tab:crossembod}
\small
\setlength{\tabcolsep}{6pt}
\begin{tabular}{lcccc}
\toprule
\textbf{Method} & \textit{pick\_place} & \textit{distribute\_bread} & \textit{sort3} & \textbf{Overall} \\
\midrule
RoboTwin & 20.00 & 6.67 & 0.00 & 8.89 \\
\textbf{CRAPE (Ours)} & 46.67 & 20.00 & 6.67 & 24.45 \\
\bottomrule
\end{tabular}
\end{table}

% ==============================================================================
\section{Simulation Performance}
\label{app:simperf}

We additionally compare code-based collection (CaP) and CRAPE purely in
simulation, on the $11$ Isaac~Lab tasks originally collected for the sim2real
pretraining stage. This isolates the quality of the generated trajectories before
any real-world transfer. Table~\ref{tab:simperf} reports per-task success rate
(smolVLA). For each task and method, we evaluate using $20$ random seeds, with
$10$ evaluation rollouts per seed.

\begin{table}[h]
\centering
\caption{Simulation success rate (\%) on the $11$ Isaac~Lab tasks (smolVLA), comparing code-based collection (CaP) and CRAPE.}
\label{tab:simperf}
\small
\setlength{\tabcolsep}{6pt}
\begin{tabular}{llcc}
\toprule
\textbf{ID} & \textbf{Task} & \textbf{CaP} & \textbf{CRAPE (Ours)} \\
\midrule
S1  & pick\_place & $10.00 \pm 5.00$ & $88 \pm 4.47$ \\
S2  & sort3 & 0 & $74 \pm 5.48$ \\
S3  & stack2 & 0 & $76 \pm 8.94$ \\
S4  & push\_button & -- & $98 \pm 4.47$ \\
S5  & turn\_on\_lever & -- & $76 \pm 11.4$ \\
S6  & turn\_off\_lever & -- & $94 \pm 5.48$ \\
S7  & extract\_block & -- & $94 \pm 8.94$ \\
S8  & open\_box & -- & $96 \pm 5.48$ \\
S9  & close\_box & -- & $66 \pm 11.4$ \\
S10 & push\_cube & -- & $100$ \\
S11 & pull\_cube & -- & $100$ \\
\midrule
\multicolumn{2}{l}{\textbf{Overall}} & -- & $87.45 \pm 6.01$ \\
\bottomrule
\end{tabular}
\end{table}

% ==============================================================================
\section{Implementation Details and Hyperparameters}
\label{app:hparams}

We report the hyperparameters in Table~\ref{tab:hparams}, using the symbols of
Section~\ref{sec:method}. Three notes: (i) we leave the ambiguity gate of
Eq.~\eqref{eq:p2-acquisition} inactive ($\tau_{\mathrm{amb}}{=}\infty$) and control
local ambiguity through the $\lambda$ term of $M_{\mathrm{MI}}$
(Eq.~\eqref{eq:p2-mi-score}); (ii) we use the DCT single step form of
$U_{\mathrm{VLA}}$ ($R{=}1$), deterministic given the DCT coefficients; and (iii)
$m_{\min}$ requires at least $1$ covered window with at least $2$ neighbors each.

\begin{table}[h]
\centering
\caption{CRAPE hyperparameters (symbols follow Section~\ref{sec:method}).}
\label{tab:hparams}
\small
\setlength{\tabcolsep}{6pt}
\begin{tabular}{lllc}
\toprule
\textbf{Stage} & \textbf{Symbol} & \textbf{Meaning} & \textbf{Value} \\
\midrule
\multirow{6}{*}{Phase 1}
 & $J$               & subgoal candidates per step                        & 64 \\
 & $K_{\mathrm{NN}}$ & kNN for state novelty $n_s$                        & 5 \\
 & $\sigma$          & subgoal perturbation width                         & 0.05\,m \\
 & --                & clip radius factor ($r_{\max}{=}$clip$\cdot\sigma$) & 1.0 \\
 & $\mathcal{D}_z$   & perturbation shape                                 & hemisphere \\
 & $L$               & skill-sequence length                              & task-derived \\
\midrule
\multirow{8}{*}{Phase 2}
 & $N$               & perturbed candidates per seed                      & 64 \\
 & $K_{\mathrm{NN}}$ & kNN for action novelty $n_A$                       & 5 \\
 & $\beta$           & action-coverage weight                             & 1.0 \\
 & $\lambda$         & ambiguity penalty weight                           & 1.0 \\
 & $\tau_{\mathrm{MI}}$  & MI admission threshold                         & 0.0 \\
 & $\tau_{\mathrm{amb}}$ & ambiguity gate (non-binding)                   & $\infty$ \\
 & $m_{\min}$        & covered windows / neighbors per window             & 1 / 2 \\
 & --                & via-offset / joint-perturb magnitude               & 0.25\,m / 0.05\,rad \\
\midrule
\multirow{2}{*}{$U_{\mathrm{VLA}}$}
 & $R$               & denoise samples (DCT single-step)                  & 1 \\
 & $h$               & action-chunk horizon                               & 50 \\
\bottomrule
\end{tabular}
\end{table}

% ==============================================================================
\section{Trajectory and State Descriptors}
\label{app:descriptors}

We summarize the descriptors in Table~\ref{tab:descriptors}. The state embedding
concatenates an $L_2$ normalized VLA prefix feature (masked mean pool of the
backbone last layer prefix tokens) with $L_2$ normalized proprioception. Each DCT
descriptor applies an orthonormal type II DCT after cubic spline resampling the
segment to length $L_0{=}50$, keeping all coefficients. The EE delta form is
translation invariant. Both distances $d$ and $d_A$ are Euclidean ($L_2$).

\begin{table}[h]
\centering
\caption{State and action descriptors. $D_{\mathrm{vla}}$ is set by the backbone checkpoint (SmolVLA).}
\label{tab:descriptors}
\small
\setlength{\tabcolsep}{6pt}
\begin{tabular}{lll}
\toprule
\textbf{Descriptor} & \textbf{Construction} & \textbf{Dim.} \\
\midrule
$\phi_{\mathrm{state}}$ & VLA prefix ($L_2$) $\oplus$ proprio ($L_2$, $5$ DoF) & $D_{\mathrm{vla}}{+}5$ \\
Joint DCT ($U_{\mathrm{VLA}}$ target) & DCT II over joint trajectory & $50{\times}6$ \\
EE delta DCT ($\psi$, MI score) & DCT II over EE pose deltas, flattened & $300$ \\
Chunk $\psi_{\xi,\tau}$ & $\phi_{\mathrm{act}}(u_{\tau:\tau+h})$, $h{=}50$ & $300$ \\
$d,\ d_A$ & Euclidean ($L_2$) & n/a \\
\bottomrule
\end{tabular}
\end{table}

% ==============================================================================
\section{Reset Time Workspace Layout Randomization}
\label{app:layout}

We randomize the initial layout at every reset so that data starts from varied
states. For each task we resample the target object placement across the reachable
workspace and keep other objects as fixed obstacles. Table~\ref{tab:layout} gives
the sampling constraints and Figure~\ref{fig:layout_random} shows the seeds for
\textit{open\_pot}.

\begin{table}[h]
\centering
\caption{Reset time layout sampling constraints.}
\label{tab:layout}
\small
\setlength{\tabcolsep}{6pt}
\begin{tabular}{lll}
\toprule
\textbf{Constraint} & \textbf{Rule} & \textbf{Value} \\
\midrule
Reachability       & inside reachable annulus                              & $[r_{\min},r_{\max}]$ \\
Edge margin        & distance from workspace boundary                      & $30$\,px \\
Obstacle avoidance & no overlap with obstacle objects                      & \texttt{is\_obstacle} \\
Seed selection     & maximize state novelty (Eq.~\eqref{eq:p1-select-subgoal}) & top scoring \\
Seeds per task     & sampled placements                                    & $N{=}20$ \\
\bottomrule
\end{tabular}
\end{table}

\begin{figure}[t]
\centering
\includegraphics[width=0.48\textwidth]{figure/seed_visualization.jpg}\hfill
\includegraphics[width=0.48\textwidth]{figure/seed_visualization_point.jpg}
\caption{Reset time layout randomization for \textit{open\_pot} ($N{=}20$ seeds,
top view). Left: bounding box view of sampled target (lid) placements. Right:
grasp point view. Cyan: reachable annulus; green: $30$\,px edge margin; red:
obstacle (pot) region.}
\label{fig:layout_random}
\end{figure}

% ==============================================================================
\section{Skill Library}
\label{app:skills}

The code generator composes task plans from a fixed library of parameterized
primitive skills (Table~\ref{tab:skills}). Each skill $z$ exposes a subgoal
parameterization and a perturbation distribution $\mathcal{D}_{z}$
(Eq.~\eqref{eq:p1-perturbed-subgoals}); the spatial skills share a common width
($\sigma{=}0.05$\,m hemisphere), though the framework permits per skill widths.

\textbf{Code generation.}
We produce task programs with an iterative vision language pipeline based on a
Gemini model: a perception stage localizes objects and grasp points in a top view
normalized frame, and a generation stage emits a Python routine that composes the
primitives of Table~\ref{tab:skills} over the detected poses, inserting an object
redetection step at each subtask boundary. The realized plan
$\pi_{\mathrm{code}}(s_0,I)=((z_1,g_1),\dots,(z_L,g_L))$ is materialized at
execution time, each call recorded as an ordinal skill $z_i$ with its $3$D subgoal
$g_i$.

\begin{table}[h]
\centering
\caption{Primitive skill library composed by the code generator. The subgoal
perturbation $\mathcal{D}_z$ is a hemisphere of radius $\sigma{=}0.05$\,m applied
to the skill's $3$-D subgoal; non-spatial skills are not perturbed (n/a).}
\label{tab:skills}
\small
\setlength{\tabcolsep}{5pt}
\begin{tabular}{lllc}
\toprule
\textbf{Skill $z$} & \textbf{Description} & \textbf{Subgoal parameters} & \textbf{$\mathcal{D}_z$} \\
\midrule
\texttt{move\_to\_initial\_state} & return to home pose                 & --                        & n/a \\
\texttt{move\_to\_position}       & move EE to xyz (opt.\ gripper)      & position $[x,y,z]$, gripper & $0.05$\,m \\
\texttt{move\_to\_pixel}          & move EE to normalized pixel         & $[y,x]\in[0,1000]$        & $0.05$\,m \\
\texttt{execute\_pick\_object}    & descend, grasp, lift                & object xyz                & $0.05$\,m \\
\texttt{execute\_place\_object}   & place, release                      & place xyz                 & $0.05$\,m \\
\texttt{execute\_place\_lid}      & seat a lid (descend, drag, release) & place xyz, pull dist.\    & $0.05$\,m \\
\texttt{execute\_press}           & torque-limited press                & position, depth           & $0.05$\,m \\
\texttt{execute\_pull}            & open drawer/door ($-x$)             & start xyz, distance       & $0.05$\,m \\
\texttt{execute\_push}            & close drawer/door ($+x$)            & start xyz, distance       & $0.05$\,m \\
\texttt{rotate\_90degree}         & in-hand $90^\circ$ rotation         & direction $\{+1,-1\}$     & n/a \\
\texttt{detect\_objects}          & runtime re-detection                & object queries            & n/a \\
\bottomrule
\end{tabular}
\end{table}

%===============================================================================

% ==============================================================================
\section{Task Visualization}
\label{app:tasks}
We performed dataset generation and model training on a suite of 11 simulation tasks and 10 real-world tasks.

\begin{figure}[t]
\centering

% Real-world task suite: shown large to match the width of the dataset table.
\includegraphics[width=\textwidth]{figure/real_tasks.png}

\vspace{0.8em}

% Simulation task suite: auxiliary visualization for the sim-pretraining tasks.
\includegraphics[width=0.95\textwidth]{figure/sim_tasks.png}

\caption{
Task visualization for the benchmark suites.
\textbf{Top:} the 10-task SO-101 real-world benchmark used for real-world
data collection, fine-tuning and policy evaluation.
\textbf{Bottom:} the 11-task Isaac Lab simulation suite used for simulation-stage
data collection, pretraining and simulation-only evaluation.
Task names are shown below each panel.
}
\label{fig:task_suites}
\vspace{-0.1in}
\end{figure}

\section{Datasets}
\label{app:datasets}

Table~\ref{tab:datasets} reports per-task statistics of the collected real-world
datasets used for fine-tuning (committed episodes, frames, and average horizon).

\begin{table}[h]
\centering
\caption{Per-task dataset statistics (committed set, $100$ episodes per task).}
\label{tab:datasets}
\small
\setlength{\tabcolsep}{5pt}
\begin{tabular}{llccc}
\toprule
\textbf{ID} & \textbf{Task} & \textbf{Episodes} & \textbf{Frames} & \textbf{Avg.\ horizon} \\
\midrule
T1  & \textit{pick\_place}      & 100 & 33,120 & 33,12s \\
T2  & \textit{distribute\_bread}& 100 & 51,788 & 51.79s \\
T3  & \textit{sort3}            & 100 & 94,144 & 94.14s \\
T4  & \textit{stack3}           & 100 & 86,402 & 86.40s \\
T5  & \textit{close\_pot}       & 100 & 42,002 & 42.00s \\
T6  & \textit{open\_drawer}     & 100 & 28,003 & 28.00s \\
T7  & \textit{prepare\_meal}    & 100 & 35,546 & 35.55s \\
T8  & \textit{fold\_towel}      & 100 & 74,094 & 74.09s \\
T9  & \textit{open\_lid}        & 100 & 33,608 & 33.61s \\
T10 & \textit{close\_drawer}    & 100 & 26,268 & 26.27s \\
\bottomrule
\end{tabular}
\end{table}

% ==============================================================================
\section{Workspace Specification}
\label{app:workspace}

\begin{figure}[t]
\centering
\includegraphics[width=\linewidth]{figure/workspace.png}
\caption{
Workspace specification. Left-1: [Intel] RealSense D435 Depth Camera. Left-2: [Innomaker] USB U20CAM-1080P. Left-3: SO101 follower robot arm. Left-4: 3D-printed workspace. Right: 3D-printing models
}
\label{fig:workspace}
\vspace{-0.1in}
\end{figure}

The benchmark uses a fixed physical workspace with 3D-printed fixtures for
repeatable object and robot placement. Table~\ref{tab:workspace} lists the
printed parts, and Figure~\ref{fig:workspace} (to be added) shows the layout,
camera placement, and calibration reference.

\begin{table}[h]
\centering
\caption{3D-printed workspace components.}
\label{tab:workspace}
\small
\setlength{\tabcolsep}{6pt}
\begin{tabular}{lll}
\toprule
\textbf{Part} & \textbf{Function} & \textbf{Material / notes} \\
\midrule
-- & -- & -- \\
-- & -- & -- \\
-- & -- & -- \\
\bottomrule
\end{tabular}
\end{table}
% TODO: add workspace layout figure (\label{fig:workspace}): camera mounts, object jigs, calibration board.

% ==============================================================================
\section{Evaluation Protocol}
\label{app:evalproto}

Each trained policy is evaluated over $30$ rollouts under matched initial
conditions, with object poses randomized within fixed ranges and an automated
reset between rollouts. Table~\ref{tab:evalproto} gives the per-task success
criterion and initial-condition randomization.

\begin{table}[h]
\centering
\caption{Per-task evaluation protocol: success criterion and initial-condition randomization ($30$ rollouts each).}
\label{tab:evalproto}
\small
\setlength{\tabcolsep}{6pt}
\begin{tabular}{llll}
\toprule
\textbf{ID} & \textbf{Task} & \textbf{Success criterion} & \textbf{Init.\ randomization} \\
\midrule
T1  & \textit{pick\_place}       & -- & -- \\
T2  & \textit{distribute\_bread} & -- & -- \\
T3  & \textit{sort3}             & -- & -- \\
T4  & \textit{stack3}            & -- & -- \\
T5  & \textit{close\_pot}        & -- & -- \\
T6  & \textit{open\_drawer}      & -- & -- \\
T7  & \textit{prepare\_meal}     & -- & -- \\
T8  & \textit{fold\_towel}       & -- & -- \\
T9  & \textit{open\_lid}         & -- & -- \\
T10 & \textit{close\_drawer}     & -- & -- \\
\bottomrule
\end{tabular}
\end{table}

\end{document}